\section{Pre-AI 与 Post-AI 的两个世界}

\subsection{新建应用 vs 遗留系统}

Birgitta 提出了一个关键的实践问题：假设我们开发出了良好的 harness 技术，能够将 AI 自主性提升到很高水平并对结果充满信心------哪些技术可以应用于\textbf{现有应用}，哪些只适用于\textbf{从零开始}、从一开始就考虑 harness 的应用？

\begin{warningbox}{遗留代码库的 Harness 改造陷阱}
对于较旧的代码库，需要慎重考虑改造（retrofit）harness 是否值得投入。虽然 AI 可以帮助加速改造过程，但这些遗留应用通常具有以下特征：
\begin{itemize}[leftmargin=1.5em]
    \item 高度非标准化的代码结构
    \item 长期积累的技术债务和熵增
    \item 缺乏文档和架构约定
    \item 隐式依赖和未记录的假设
\end{itemize}
这就像在一个从未使用过静态代码分析工具的代码库上首次运行分析------你会被大量告警淹没。
\end{warningbox}

\subsection{两个世界的分化}

这一观察暗示了未来可能出现\textbf{两种截然不同的软件维护模式}：

\begin{table}[H]
\centering
\begin{tabular}{lll}
\toprule
\textbf{维度} & \textbf{Pre-AI 应用} & \textbf{Post-AI 应用} \\
\midrule
代码生成 & 主要由人类编写 & 主要由 AI agent 生成 \\
架构约束 & 隐式、基于约定 & 显式、由 harness 强制 \\
维护方式 & 人工审查 + 重构 & agent 驱动 + 自动清理 \\
Harness 适配 & 需要昂贵的改造 & 从设计之初内建 \\
技术债务 & 大量历史遗留 & 由 ``垃圾回收'' agent 持续清理 \\
\bottomrule
\end{tabular}
\end{table}

对于组织而言，这意味着需要制定\textbf{差异化的策略}：对新项目采用 harness-first 的方法，对遗留系统则需要评估改造的 ROI（投资回报率）。

\subsection{本章小结}

harness 技术在新建应用和遗留系统上的适用性存在本质差异。组织可能面临``两个世界''的分化：harness-native 的新应用和难以改造的旧应用。对遗留系统的 harness 改造需要谨慎评估投入产出比。

\newpage

